\subsection{From document provenance to epistemic provenance}

Machine-readable provenance solves an important but limited problem: it records where an object came from and how it was transformed. It does not, by itself, determine what the object is allowed to support scientifically. The Open Research Knowledge Graph (ORKG), for example, moves scholarly communication toward structured, machine-actionable research contributions rather than document-only representation \cite{jaradeh2019orkg}. W3C PROV supplies interoperable entity--activity--agent derivations \cite{w3cprov2013}. Nanopublications and RO-Crate provide other granular or packaged representations \cite{kuhn2018,soilandreyes2022}. These are essential precedents for a scientific state that must be inspectable by software.

\RAKL{} adds a second layer that we call \emph{epistemic provenance}. For an asserted scientific relation, the system must retain not only the source chain but the context alignment, verification identity, relation type, error semantics and authority effect. A source can be perfectly provenance-traceable and still be weak evidence; two sources can be independently published but descend from the same primary dataset; a derived analysis can be reproducible while its measurement operator is mismatched to the claim. Thus
\[
\text{traceable provenance}\not\Rightarrow
\text{independent evidence}\not\Rightarrow
\text{sufficient authority}.
\]
The evidence-lineage graph addresses the middle implication. It attempts to represent shared ancestry so that ten derivative reports of one measurement do not count as ten independent flat searches or ten independent confirmations.

Cognitive provenance is separated again. A workspace item may causally influence a proposal because it was retrieved or emphasized, yet that influence does not make the item evidence for the proposal's scientific content. The two edges have different semantics:
\[
\Pi^{\mathrm{cog}}\neq \Pi^{\mathrm{evid}}.
\]
This distinction becomes important when using LLM traces, memory systems or interpretability tools. A representation being computationally load-bearing is evidence about the agent's cognition; it is not automatically evidence about the external phenomenon being studied. A paper that mixes those graphs can accidentally turn explanation of the model into explanation of nature.

The practical consequence is an ``authority firewall'' around transformations. Raw acquisition can create a source object; projection can create a context-scoped claim object; normalization can create an equivalence witness; an embedding can create a retrieval index; compression can create a derivative view; a workspace can create an access relation. None of those operations, on its own, increases scientific authority. Only a verified promotion operation can do so. This makes provenance a necessary substrate for authority without making provenance itself authority.

\subsection{Analogy and active inquiry}
Structure-mapping theory treats analogy as relational mapping rather than surface resemblance \cite{gentner1983}; constraint-satisfaction accounts jointly evaluate competing mappings \cite{holyoak1989}. This prior art motivates RAKL's separation of conservative `GLUE' from exploratory `JUMP': an analogy must state preserved and non-preserved structure and has discovery authority only until target evidence arrives. Experimental information itself predates modern active learning: Lindley formalized expected information from an experiment \cite{lindley1956}, and Bayesian experimental-design reviews systematized utility-based design under prior/posterior uncertainty \cite{chaloner1995}. Active data selection also has a long history: information-based design chooses observations by their expected information value relative to a model class \cite{mackay1992}. RAKL uses information gain when calibrated probabilistic models are defensible, and otherwise falls back to set shrinkage or worst-case mechanism separation rather than fabricating priors. Repeated adaptive evaluation is itself a validity problem, motivating the protected fresh-assurance reserve \cite{dwork2015}.

\subsection{Discovery search, experimental design and stopping are one control problem at different layers}

Literature search and physical experimentation are often treated as separate modules, but both can be written as selection under uncertainty. A literature query chooses an observation operator over the published record; an experiment chooses an observation operator over the scientific system. In each case the purpose is not simply to collect more items. It is to expose distinctions that matter for the registered target. Bayesian experimental design formalizes one version through expected utility \cite{lindley1956,chaloner1995}; active learning chooses measurements that are expected to improve a learner \cite{mackay1992}. \RAKL{} uses these ideas only when the probability model required by the utility is itself defensible.

Literature-based discovery supplies an important cross-domain precedent for the search side. Swanson's classic ``undiscovered public knowledge'' example showed that useful connections can exist between literatures that are individually public yet not directly connected \cite{swanson1986}. This is precisely why semantic saturation cannot be established by repeating the framework's own vocabulary. Open-World Mechanism Discovery therefore includes function-only, historical, mathematical, implementation, citation-neighborhood, adversarial and freshness routes. A later exogenous concept that changes the novelty boundary is a falsifier of the prior route-closure claim for that scope.

Stopping research introduces a dual risk. Stopping too early misses a material semantic object; stopping too late wastes resources and creates more opportunities for adaptive overfitting. Technology-assisted systematic review has developed explicit stopping methods because screening ``until tired'' is not a defensible rule. Recent confidence-based approaches shift attention from raw recall toward whether the information need is sufficiently resolved \cite{fletcher2026stopping}. \RAKL{} adopts the general lesson without importing a particular recall estimator: stopping must be attached to a registered information need, source/query universe, evidence lineage and reopen condition.

This produces a hierarchy of stopping statements. ``Budget exhausted'' is a resource state. ``Route locally flat'' means one query family yielded no new canonical semantic object. ``Same-context plateau'' means repeated passes in one research context are flat. ``Evidence-lineage independent flatness'' requires genuinely distinct evidence ancestry. ``Scoped saturation'' additionally requires route coverage, unresolved-obligation accounting and no native residual that reopens the fiber. None of these implies that the open world contains no unknown relevant concept.

\subsection{Metacognition and learning from challenge}
RAKL's human-learning analogy is functional rather than anthropomorphic. Self-regulated learning distinguishes forethought, monitoring/control and reflection \cite{zimmerman2002}; productive-failure work shows that initial failure can improve later representation and transfer under appropriate conditions \cite{kapur2008}. Psychology also documents reasons not to trust verbal introspection alone: people can underestimate their own bias \cite{pronin2002}, overestimate explanatory understanding until forced to reconstruct an explanation \cite{rozenblit2002}, and benefit in some settings from explicitly considering the opposite hypothesis rather than generic instructions to ``be unbiased'' \cite{lord1984}. These results motivate executable RAKL controls such as calibration audits, explanation reconstruction, countermodels, strategy switching, help-seeking, anti-rumination and fresh transfer. Recent generative-AI work similarly argues for metacognitive control around rather than merely inside the model \cite{ji2026}.

\subsection{Why reflection is not an evaluator}

A self-reflective LLM can notice errors, but the act of reflection is not statistically or epistemically independent of the process that generated those errors. This matters because many contemporary agent loops reuse one model as proposer, critic, judge and editor. Such role separation can be operationally useful while leaving shared blind spots, shared training priors and shared context contamination intact. The psychology literature on bias blind spots and illusion of explanatory depth provides human analogies for this failure \cite{pronin2002,rozenblit2002}; \RAKL{} treats the analogy as motivation, not as a claim that language models reproduce the same cognitive mechanisms.

The architecture therefore separates \emph{metacognitive diagnosis} from \emph{assurance}. Diagnosis asks what kind of failure occurred and what action should be taken next. Assurance asks whether a frozen candidate method change transfers on a task or realization that did not participate in designing the repair. The first can be same-context. The second must be protected to the degree required by the claim. Development improvement \(\Delta_D>0\) is thus necessary but insufficient for strong self-evolution; the protected comparison must also show \(\Delta_A>0\) under unchanged blocking criteria.

This separation also prevents ``reflection inflation.'' Adding more self-critique turns is not automatically more metacognition. A useful metacognitive operation must change a control decision: acquire evidence, switch strategy, construct a countermodel, ask for external help, stop an unproductive route, or open a method-basis gap. Repeating explanations without changing an uncertainty-bearing operation is recorded as flat effort. That is why learning-progress signals are coupled to action policy rather than treated as authority.

Finally, metacognitive state itself must remain revisable. A controller can misattribute a data problem to a method problem or a method problem to a model problem. Self-\RAKL{} therefore treats failure attribution as a hypothesis with discriminators. Only a supported method-basis gap justifies inventing or assimilating a new operator. Missing measurements route to evidence acquisition; implementation defects route to repair; objective defects route to governance. The ability to invent a method is not permission to invent one whenever performance is poor.

\subsection{Memory, retrieval and local-to-global knowledge}
Retrieval-augmented generation already combines parametric generation with explicit non-parametric memory \cite{lewis2020}; therefore external knowledge access is not itself a RAKL contribution. Hierarchical memory and context optimization are prior engineering domains. MemGPT frames long-horizon context as virtual memory \cite{packer2023}; RAPTOR organizes multi-resolution retrieval hierarchically \cite{sarthi2024}; LLMLingua compresses prompts \cite{jiang2023}; and RECOMP selectively compresses retrieved material \cite{xu2023}. Long-context models also do not necessarily use all prompt positions equally: ``lost in the middle'' evaluations show systematic degradation when relevant information lies away from context boundaries \cite{liu2024}. This supports treating active scientific context as an engineered working set rather than assuming that a maximal context window is equivalent to usable memory. Sheaf-theoretic work has long characterized contextuality through obstructions to global sections \cite{abramsky2011}; Local-to-global consistency and provenance-rich contextual worlds are also active formal directions \cite{gibson2026,vitali2026}. RAKL therefore does not claim generic graph storage, compression, hierarchical retrieval or sheaf mathematics. Its narrower requirement is that scientific projection, semantic identity, storage compression and prompt selection remain distinct operations, and that mandatory epistemic evidence cannot be silently traded away for context-window efficiency.

Taken together, these literatures narrow the candidate novelty to the integrated transition discipline: source contribution as contextual projection before competition; typed non-escalating authority; obstruction-aware local-to-global synthesis; immutable negative research history; evidence-lineage-aware stopping; bounded but rehydratable scientific memory; residual-driven method completeness; and protected recursive method evolution.

\paragraph{Compression lineage and what is not claimed.}
Minimum-description-length reasoning formalizes an important pressure against representational complexity \cite{rissanen1978}, while the information-bottleneck principle studies compression that preserves information relevant to a target variable \cite{tishby2000}. These are intellectual-lineage analogies rather than algorithms silently implemented by the current reference profile. The present context compiler does not estimate mutual information or solve an information-bottleneck objective, and the Knowledge Atlas ``volume'' introduced below is not a description length. The retained engineering lesson is narrower: compactness, target-relevant retained function and scientific authority must be measured separately.
